\chapter{综合结论与后续验证}
\section{已形成的控制与验证链路}
\subsection{轨迹与反馈系统}
ROS2控制节点与Isaac物理仿真之间的目标、状态和时钟接口已经建立。针对初期跟踪滞后，参考速度与驱动增益调整在独立运动幅度下将平均运动误差从0.864度降至0.082度。该结果限定在同一Panda模型与场景，不能将当前驱动参数直接用于另一台机械臂。

状态反馈部分连接了虚拟IMU、公开真实数据和Jetson实体采集。最明确的工程结论是信息新鲜度与发布频率不同：实体传感器约25赫兹输出，200赫兹发布的绝大多数消息复用测量；五点平均又会引入约80毫秒群延迟。消息传输不足1毫秒的第95百分位，并不意味着整条测量到执行链路不足1毫秒。
\subsection{接触控制}
法向导纳能降低持续压入力，但空中直接启用力误差积分会加快下压，反而增大首次冲击。接触门控、接近减速、目标力缓升、连续逆解和重力补偿分别处理不同环节，不宜归结为一个参数的作用。

后期选用条件在1毫秒步长下原始峰值11.11牛、末段均力约5.007牛，满足项目预设的四项检查；这不是工业认证标准。刚性接触减半步长后峰值达到20.43牛，部分高度与摩擦条件也暴露敏感性。柔顺材质改善了这组数值稳定性，但改变环境不能算作刚性控制器本身已经解决全部问题。
\subsection{从简化平衡到实际多关节行走}
线性倒立摆提供了受约束的支撑点与落脚分析，有限脚掌模型补充了实际接触。随后六电机平面双足将摆脚动力学、全身逆动力学和真实脚地接触连通，八步标准条件前进约0.937米；后期两组二十步均前进约2.323米。

显式零力矩点任务在100赫兹无扰动条件下将实际压力中心误差由19.41毫米降至6.06毫米，但质心误差升高，后向恢复也更慢。94个去重条件79个通过，长延迟仍频繁失败。这说明物理量接口统一后，仍需要面对多任务权重、接触切换与状态预测之间的折中。
\subsection{三维全身控制}
G1模型保留自由浮动基座、实际碰撞和29个受限电机，完成左右单脚承重和三维慢速交替迈步。四步参考29秒前进约0.210米，六步43秒约0.330米。这些结果超出了平面站立，但仍属于准静态步行，没有验证动态高速运动。

正式52个条件通过51个，另8个强扰动仅通过2个。求解器残差门槛、自接触和支撑时序都是实际限制。Clarabel在所测条件的调用耗时尾部较短，但不包含全部传感、通信与调度，不能据此称为硬实时控制器。
\section{跨阶段得到的工程经验}
\begin{enumerate}
\item 先核对单位、符号与执行通道。地面支撑点的单位为米，接触扳手与电机力矩也不是同一变量；只有经过完整动力学约束才可转换成执行指令。
\item 按假设逐层增加真实性。关闭重力、简化碰撞和限制平面自由度有助于隔离问题，但恢复这些条件后必须重新验证。
\item 平均指标与尾部风险同时保留。原始力峰值、求解耗时尾部、提前终止和接触确认不能被均值掩盖。
\item 控制器只能使用当时可获得的信息。延迟队列、旧帧复用、漂移和预测模型的输入应明确，不能暗中回读真实当前状态。
\item 将数值核验与任务验收分开。接触力重算一致和动力学残差很小，说明实现自洽；是否走完、是否恢复承重与是否满足时延仍须各自判断。
\end{enumerate}
\section{仍需补齐的部署条件}
\subsection{状态估计与接触事件}
将IMU、关节编码器与足部接触接入同一状态估计器，明确漂移、偏置和丢包；以实际触地事件对名义时序做有限调整。现有测量测试和控制仿真分别开展，还没有完成这条真实传感闭环。
\subsection{执行器与额外接触}
有限力矩不等于真实电机模型。下一阶段应加入带宽、通信、速率和热约束，并对自碰撞、提前触地和地形变化建立控制约束。强扰动中已观察到的腕部与髋部、肩部与躯干接触提供了具体检查对象，但尚不能认定为所有失败的唯一原因。
\subsection{接触控制的独立验证}
优先验证未知表面高度、斜面、切向运动、失去接触后的退出与恢复，并用已标定的有限刚度及测量带宽检验力峰值。固定平面单法向结果不自动推广到六维柔顺操作。
\subsection{真机测试的准备}
如获得合适设备，先确认机器人的力矩接口、内置重力补偿、状态时间戳和停止方式，再选择低风险自由空间轨迹与有限接触工况。测试材料应包含接口说明、参数限制、数据记录和明确验收条件。当前未获得代测结果，也没有把方案准备计为真实硬件通过。
\section{报告与代码的使用方式}
各阶段的相关代码和源报告按附录A索引；模型来源、配置和依赖保留原位置。附录B列出已复制图片与英文命名视频，附录C解释术语，附录D保留历史运行示例。源码目录不等于包含所有资产和环境的独立软件发行包，复现应先阅读对应README，并为新运行建立新输出目录。

\section{官方行走基线补充结论}
官方强化学习策略的动态迈步已复现；0.5米/秒工况20秒前进9.23米、横移1.41米，航向变化约15.6度。归零指令后平均净移动减慢，但仍保持踏步，未达到预设瞬时停止条件。自编全身控制保持准静态原型定位；本阶段收尾后，不继续扩大自然步态开发，保留模式切换、方向保持、状态估计和真机验证作为明确边界。
